\chapter{绪论}

% 官方要求：明确技术方案的背景和目标，围绕两个问题：为什么要做，以及做到什么程度（目标）。

\section{赛季背景：空中机器人定位与规则的演变}

RoboMaster 2026 赛季空中机器人相关规则较上赛季国赛发生了较大改动：取消了冷却时间限制，引入了与地面兵种一致的热量与经验等级体系；空中支援由"限次"改为"限时"，可通过免费时间或金币购买时间直接召唤；17mm 小弹丸兵种的基础伤害由 10 提升至 20，但基地装甲板由大装甲板改为小装甲板，且基地顶部装甲板受小弹丸伤害降为一发 5 血。基础伤害的大幅提升显著增强了空中机器人对前哨站与地面兵种的威胁，但也使无人机"磨基地血量"的收益大幅下降；热量冷却与经验等级的引入，则要求空中机器人具备更高的持续作战与精准打击能力。

此外，本赛季新增了雷达反制无人机机制，限制了空中机器人云台长时间对地面的探测与火力压制，对无人机的机动性与结构稳定性提出了更高要求；飞手耳机的取消，要求飞手与云台手共同建立一套可靠的赛场交流与战术协同体系。

在赛场地形方面，本赛季新增了 \qty{70}{\mm} 高的起伏路段，隧道高度降至 \qty{250}{\mm}，限制了地面兵种的大范围机动，使空中机器人的打击目标移动范围更加集中、更易击杀地面单位，也进一步凸显了空中机器人对自瞄与弹道精度的要求。

综合上述规则与地形变化，空中机器人本赛季已从"辅助单位"向"核心战力"转型，成为集视野侦察、对地火力压制、前哨站攻坚、能量机关开启与基地打击于一体的多功能作战兵种。相比十字型布局，本赛季主流的四轴 X 型机架布局具有更大的桨盘流通面积、更优的姿态控制裕度与更好的云台视野，是仲恺奇点、深技大悍匠等强队以及本队空中机器人共同采用的构型方案。

\section{空中机器人的研发背景与目标}

\subsection{战术定位}

结合赛季规则与场地地形，本赛季空中机器人的主要战术价值体现在：
\begin{itemize}
    \item \textbf{开局压制}：比赛前 2 分钟优先打击敌方工程机器人等高价值地面单位，重创敌方经济与经验体系，为我方地面部队打开局面；
    \item \textbf{前哨站攻坚}：对敌方前哨站进行补刀或摧毁，为飞镖、英雄等进攻型单位创造机会；
    \item \textbf{对地火力压制}：支援地面团战、补刀残血敌方地面单位，防守时骚扰敌方吊射英雄、保护我方关键单位；
    \item \textbf{能量机关与基地}：凭借大角度双轴云台，承担开启能量机关以及对敌方基地的打击任务。
\end{itemize}

为此，空中机器人必须具备良好的悬停稳定性与续航能力、精准的弹道与自瞄，以及高可靠的结构与定位系统，才能在场上持续发挥作用。

\subsection{分级设计目标}

针对上述定位，我们对本赛季空中机器人制定了分级设计目标，如\tabref{tab:绪论分级设计目标}所示。

\begin{longtable}[htbp!]{ m{2.6cm}<{\centering}  m{2.6cm}<{\centering}  m{9.6cm}<{\centering} }
\hline
需求ID & 需求 & 需求目标 \\
\hline
\endfirsthead
\hline
需求ID & 需求 & 需求目标 \\
\hline
\endhead
\hline
\multicolumn{3}{r}{接下页}
\endfoot
\caption{空中机器人分级设计目标}
\label{tab:绪论分级设计目标}
\endlastfoot
 \hline
 机组-001 & 强续航 & 实战场景下飞行时间达到 7 分钟以上，动力系统稳定、抗干扰能力优秀 \\
 \hline
 机组-002 & 定点定位 & 搭载 GPS、Livox MID-360 雷达与激光检测模块，感知自身位姿，保持定点稳定飞行 \\
 \hline
 机组-003 & 轻量化高强度 & 以碳纤维板为主体结构，整体刚性强、质量轻，满足检录收纳与快速部署要求 \\
 \hline
 机组-004 & 折叠收纳 & 机臂与桨保可折叠，折叠后尺寸满足检录最大收纳尺寸要求，折叠操作简单 \\
 \hline
 机组-005 & 弹道精度 & 高频顺畅发射不卡弹，\qty{9}{\m} 距离散度控制在小装甲板范围内 \\
 \hline
 机组-006 & 大角度云台 & Yaw 轴水平 $180^{\circ}$ 旋转，Pitch 轴俯仰角满足近距离开启能量机关的需求 \\
 \hline
 机组-007 & 安全可靠 & 建立"定位—指令—权限"三级安全保护与一键夺权机制，实现"安全飞行，永不坠机" \\
 \hline
\end{longtable}

\section{目标实现计划}

围绕上述分级目标，本赛季空中机器人的主要实现计划如下：
\begin{itemize}
    \item 完成可折叠机臂与分体式桨保设计，在保证结构强度的前提下实现轻量化，满足检录收纳尺寸与快速部署要求；
    \item 基于 FAST-LIO 激光惯性里程计与高频 IMU 插值，实现全场连续、可靠的定位与状态估计，保障定点悬停与航点飞行；
    \item 完成横向双相机大角度双轴云台与高精度发射系统设计，实现 \qty{9}{\m} 距离小装甲板散度内的高频、高精度打击；
    \item 完成自瞄系统（装甲板、前哨站、能量机关等多任务链）与飞行指令桥接，实现手动、自瞄、自动飞行等模式的平滑切换；
    \item 建立"定位—指令—权限"三级安全保护与遥控器一键夺权机制，从算法与软件层面杜绝失控坠机风险。
\end{itemize}
